\section*{Chapter \ref{ch:in:career-paths-and-moves} --- Career Paths and Moves}

\begin{solution}{exo:in:career-paths-and-moves:1}
$0.4\times350\,000\times(3+1)/2=\$280\,000$ unvested. Resigning at the start of March, after the February payment,
forfeits those and two months of accrual, $350\,000\times2/12=\$58\,333$: \$338\,333.
\end{solution}

\begin{solution}{exo:in:career-paths-and-moves:2}
Seven weeks: one week for each year of service between two and twelve years.
\end{solution}

\begin{solution}{exo:in:career-paths-and-moves:3}
Twelve months without an award: \$350\,000.
\end{solution}

\begin{solution}{exo:in:career-paths-and-moves:4}
The buyout replaces the forfeited deferrals and the unpaid award (\$630\,000), but not the month's accrual
(\$29\,167) or the year out of the market (\$350\,000): the February move still costs \$379\,167, a 21.1\% raise,
against 38.2\% for a March move without a buyout.
\end{solution}

\begin{solution}{exo:in:career-paths-and-moves:5}
In none: with a 50\% buyout chance the neutral raise is at least 30.5\% (March) and up to 45.0\% (December).
\end{solution}

\begin{solution}{exo:in:career-paths-and-moves:6}
An unpaid non-compete adds the base for its months to the cost (\$125\,000 for six months here); a paid one costs only
the award that does not accrue. Some legal systems also look at payment when deciding whether a clause is reasonable.
\end{solution}

\begin{solution}{exo:in:career-paths-and-moves:7}
The cost rises by \$125\,000 to \$819\,167, and the neutral raise to 45.5\%.
\end{solution}

\begin{solution}{exo:in:career-paths-and-moves:8}
The chain's transition probabilities are the chapter's assumptions, not estimates: 21.5\% is what those assumptions
imply, and says nothing about real careers. The model is useful for the shape of the fan and for sensitivity to the
assumptions.
\end{solution}

\begin{solution}{pb:in:career-paths-and-moves:1}
\begin{enumerate}
\item As in the chapter's definition.
\item As trials for full-time roles, and graduate roles built around learning through real work, in several offices.
\item To research roles, above a new graduate's level.
\item Bank to buy side and back, between buy-side firms, between technology and trading.
\item None; only the cost of a move can be measured.
\item As in the definition; one week a year of service up to twelve weeks from the employer, one week from the employee.
\item As in the method.
\item February: \$659\,167 forfeited, \$694\,167 net with a 50\% buyout chance. March: \$338\,333 forfeited, \$548\,333
net.
\item 38.6\% and 30.5\%.
\item The year out of the market, \$350\,000 of award, because it is lost in every month and no buyout replaces it.
\item Stopped in August 2024 and vacated; the FTC dismissed its appeals in September 2025.
\item A three-month limit announced in May 2023 and not enacted; a working paper of November 2025 on limits, a ban or a
salary threshold.
\item Such clauses are unenforceable in California wherever they were signed, and trying to enforce them is a civil
violation.
\item About nine months (0.74 years), from the strategy's profit, the share lost, the edge's half-life and the pay.
\item Capital, investors, a track record they believe, and the freedom from restrictive covenants to use what one knows.
\item Between 21.1\% and 56.1\% at the start of February depending on the buyout; 38.6\% at a 50\% chance, against 30.5\%
a month later.
\item The forfeited deferrals and the unpaid award, and ideally the accrual, in cash or in the new firm's deferred pay.
\item It would keep the deferrals vesting, removing \$280\,000 from the cost.
\item The shape of a career under stated assumptions: a widening fan and a large share of exits; not real frequencies.
\item Wait until the award is paid unless the new firm buys out all of it, and ask for the accrual too. The year out of the
market is the largest cost, so negotiate its length and whether it is paid.
\end{enumerate}
\end{solution}

\begin{solution}{iq:in:career-paths-and-moves:1}
Say what you want to do that the new firm offers, not what is wrong with the old one; be specific and consistent.
\iqlookfor{a forward-looking reason.}
\end{solution}

\begin{solution}{iq:in:career-paths-and-moves:2}
The kind of work, the markets and horizons, the methods in general terms and your role; never the signals, parameters,
code, data or positions, which belong to your employer.
\iqlookfor{respect for confidentiality.}
\end{solution}

\begin{solution}{iq:in:career-paths-and-moves:3}
Study, rest and prepare in ways the clause allows: courses, public research, tools unrelated to the old employer's
strategies; not trading or working for a competitor.
\iqlookfor{knowing what the clause forbids.}
\end{solution}

\begin{solution}{iq:in:career-paths-and-moves:4}
With the employer's permission or with aggregates that reveal nothing proprietary (Sharpe ratio, drawdown, capacity,
turnover), references, and a description of the process; never with the employer's data.
\iqlookfor{evidence without disclosure.}
\end{solution}

\begin{solution}{iq:in:career-paths-and-moves:5}
Closer ownership of results and faster feedback; you would miss the breadth of products, the client flow and the bank's
infrastructure.
\iqlookfor{an honest trade-off.}
\end{solution}

\begin{solution}{iq:in:career-paths-and-moves:6}
List the unvested tranches with their vesting dates and current value, the unpaid award, and the accrual; propose a buyout
in the new firm's deferred pay on a matching schedule, so that both sides keep the incentive.
\iqlookfor{a concrete schedule and aligned incentives.}
\end{solution}
